% \blindtext
\addcontentsline{toc}{subsection}{Theorem name}
\begin{thm}[Eindeutigkeitssatz zu $l^p$ Metriken]\label{lp-eindeutig}
Sei $\langle A, \succsim \rangle$ eine vollständige, segmentally additive proximity structure, die auch eine additive difference structure ist.\\
Ferner sei $\delta$ homogen, also $\delta(x, x + t(z-x)) = t \delta(x,z)$, für $t > 0$.\\
Dann gibt es ein eindeutiges $r \geq 1$ und reellwertige Funktionen $\varphi_i$, die auf $A$ definiert sind, $i = 1, \ldots, n$, die bis auf eine positive lineare Transformation eindeutig sind, so dass für alle $a, b, c, d \in A$,
\[
(a, b) \succsim (c, d) \iff \delta(a, b) \geq \delta(c, d),
\]
wobei
\[
\delta(a, b) = \left[\sum_{i=1}^{n} |\varphi_i(a) - \varphi_i(b)|^r\right]^{1/r}.
\]
\end{thm}

% \blindtext

\begin{rem}
  Dieser Satz ist ursprünglich ohne die Zusatzannahme der Homogenität von $\delta$ formuliert. Jedoch konnten wir den Beweis an einer Stelle nicht nachvollziehen bzw. denken dass er möglicherweise sogar falsch ist, und nehmen daher zusätzlich diese Eigenschaft an.
\end{rem}

\begin{proof}
  Der Beweis passiert in mehreren Schritten. Zuerst sagen wir etwas über die Definitions- und Zielbereiche der Funktionen $F_i$ und finden damit Grenzen für den Definitionsbereich von $F$.
  Danach argumentieren wir, dass wir $F$ bis auf die positiven reellen Zahlen $[0, \infty)$ erweitern können und dass diese Erweiterung eine Funktionsgleichung impliziert, die nur von $f(x) = kx^p$ erfüllt wird.
  Im letzten Schritt zeigen wir, dass $F$ auf ganz $\mathbb{R}^+$ gleich der Erweiterung sein muss und damit $\delta$ eine $l^p$ Metrik sein muss.

Nach Satz \todo{ref hinr-notw-beds} liefert das additive difference Model

\[
\delta(a, b) = G \left( \sum_{i=1}^{n} F_i \left( |\varphi_i(a_i) - \varphi_i(b_i)| \right) \right)
\]

nur dann eine Metrik, wenn $F_i(\alpha) = F(t_i \alpha)$ und $F = G^{-1}$ gilt. Daher ist $F^{-1}$ für alle Zahlen $\sum_{i=1}^{n} F_i \left( |\varphi_i(a_i) - \varphi_i(b_i)| \right)$ definiert, und jedes $F_i$ stimmt mit $F$ auf seinen Definitionsbereich überein (bis auf den Faktor $t_i$, der das Argument multipliziert und den Definitionsbereich anpasst).\\
Da die Definitionsbereiche von $F_i$ Intervalle sind (weil $\delta$ eine Metrik mit additiven Segmenten ist), können wir $\omega > 0$ wählen, so dass für jedes $i = 1, \ldots, n$ und jedes $\alpha$ mit $0 \leq \alpha \leq \omega/t_i$ gilt, dass $\alpha$ im Bereich von $F_i$ liegt.\\
Insbesondere, wenn wir $\omega_i = \omega/t_i$ setzen, ist $F^{-1} \left[ \sum_{i=1}^{n} F_i(\omega_i) \right] = F^{-1} \left[ nF(\omega) \right] =: \Omega$ definiert. Damit ist $F = G^{-1}$ mindestens im Intervall $[0, \Omega]$ definiert. Wir können $F$ nun wie folgt erweitern.

Für $\alpha < \omega$ gilt mit der Homogenität

\begin{align*}
  F^{-1}[nF(\alpha)] &= F^{-1}[nF(\alpha \omega/\omega)] \\
                     &= (\alpha/\omega) F^{-1}[nF(\omega)] \\
                     &= \alpha \Omega/\omega.
\end{align*}

Das liefert durch Einsetzen der korrekten Werte für $\alpha$ ($\tilde  \alpha = \alpha \frac{\omega}{\Omega}$ und $\tilde \alpha = F^{-1}( \frac{\alpha}{n} )$)


\begin{equation}
  F(a) = nF(a \omega/\Omega), \quad 0 \leq a \leq \Omega; \label{f-erw}
\end{equation}
\begin{equation}
  F^{-1}(a) = (\Omega/\omega) F^{-1}(a/n), \quad 0 \leq a \leq nF(\omega). \label{f-1-erw}
\end{equation}

Diese beiden Gleichungen können wir jetzt nutzen, um den Definitionsbereich von $F$ und $F^{-1}$ zu erweitern.
Wir nutzen (\todo{ref f-erw}) indem wir $F$ durch die rechte Seite definieren (für Werte, bei denen die rechte Seite definiert ist, also für $0 \leq a \leq \Omega^2/\omega$). Da $\Omega^2/\omega > \Omega$ gilt erweitert dies $F$ wirklich über das Intervall $[0, \Omega]$ hinaus.\\
Entsprechend erfüllt $F^{-1}$ Gleichung (\todo{ref f-1-erw}) für das erweiterte Intervall $0 \leq a \leq nF(\Omega) = n^2 F(\omega)$.\\
Darüber hinaus gilt die Homogenität nun für $0 \leq t \alpha_i \leq \Omega$:

\begin{align*}
  F^{-1} \left[ \sum_{i=1}^{n} F(t \alpha_i) \right] &= F^{-1} \left[ \sum_{i=1}^{n} nF(t \alpha_i \omega \backslash \Omega) \right] \quad (\text{nach Gleichung \todo{ref f-erw}}) \\
                                      &= F^{-1} \left[ n \sum_{i=1}^{n} F(t \alpha_i \omega \backslash \Omega) \right] \\
                                      &= \frac{\Omega}{\omega} F^{-1} \left[ \sum_{i=1}^{n} F(t \alpha_i \omega \backslash \Omega) \right] \quad (\text{nach Gleichung \todo{ref f-1-erw}}) \\
                                      &= \frac{\Omega}{\omega} t F^{-1} \left[ \sum_{i=1}^{n} F(\alpha_i \omega \backslash \Omega) \right] \quad (\text{wegen Homogenität, da } \alpha_i \omega \backslash \Omega \leq \omega) \\
                                      &= t F^{-1} \left[ n \sum_{i=1}^{n} F(\alpha_i \omega \backslash \Omega) \right] \quad (\text{nach Gleichung \todo{ref f-1-erw}}) \\
                                      &= t F^{-1} \left[ \sum_{i=1}^{n} n F(\alpha_i \omega \backslash \Omega) \right] \\
                                      &= t F^{-1} \left[ \sum_{i=1}^{n} F(\alpha_i) \right] \quad (\text{nach Gleichung \todo{ref f-erw}}) \\
\end{align*}


Durch wiederholte Anwendung erweitern sich die Intervalle, in denen $F$ definiert ist und Homogenität gültig ist, jedes Mal um einen Faktor $\Omega/\omega$. Somit können wir $F$ auf $[0, \infty)$ erweitern sodass die Homogenität erhalten bleibt.

Die Funktion, die durch das erweitern entstanden ist, muss somit mindestens auf dem Intervall $[o,\infty)$ mit der ursprünglichen Funktion $F$ übereinstimmen.
Außerdem gilt mit Gleichung \todo{ref f-1-erw} durch Einsetzen der Mittelwertsfunktion $F^{-1} \left[ \left( \frac{1}{n} \sum_{i=1}^{n} F(\alpha_i) \right) \right]$:

$$
  F^{-1} \left[ \frac{1}{n} \sum_{i=1}^{n} F(t \alpha_i) \right] = tF^{-1} \left[ \frac{1}{n} \sum_{i=1}^{n} F(\alpha_i) \right].
  $$

  Nach \cite[Satz 84]{hardy1952inequalities} sind die einzigen monotonen Funktionen, die diese Gleichung für $0 \leq \alpha < \infty$ erfüllen, Funktionen der Form $k\alpha^r$ (siehe auch \cite[S. 153]{aczel1966lectures}).

  Damit gilt, dass $F(\alpha) = k \alpha^r$ mindestens auf dem Intervall $[0, \Omega]$ gelten muss.
Aber jetzt können wir, für jedes $|\varphi_i(a_i) - \varphi_i(b_i)|$, ein $t > 0$ so klein finden, dass $t|\varphi_i(a_i) - \varphi_i(b_i)| < \omega, i = 1, \ldots, n$. Anwenden von der Monotonie liefert (wir können annehmen, dass $t \leq 1$)

\begin{align*}
  \delta(a, b) &= F^{-1} \left( \sum_{i=1}^{n} F \left( |\varphi_i(a_i) - \varphi_i(b_i)| \right) \right) \\
               &= \frac{1}{t} F^{-1} \left( \sum_{i=1}^{n} F \left( t |\varphi_i(a_i) - \varphi_i(b_i)| \right) \right) \\
               &= \frac{1}{t} \left( \sum_{i=1}^{n} \left( t|\varphi_i(a_i) - \varphi_i(b_i)| \right)^r \right)^{1/r} \quad \text{weil } t |\varphi_i(\alpha_i) - \varphi_i(\beta_i)| \leq \omega\\
               &= \left( \sum_{i=1}^{n} |\varphi_i(a_i) - \varphi_i(b_i)|^r \right)^{1/r} \\
\end{align*}

Damit ist $\delta$ eine $l^p$ Metrik, und durch Satz \todo{ref add-diff-prox-struct-repr} ist $F= k \cdot \alpha^r$ und damit $F_i = k (t_i \alpha)^r$ oder $F_i(\alpha) = k \alpha^r$, wenn wir die $t_i$ in die $\varphi_i$ einbauen. Schließlich folgt die Einschränkung $1 \leq r < \infty$ aus der Dreiecksungleichung oder, genauer gesagt, aus der Minkowski-Ungleichung.

\end{proof}

